\clearpage\begingroup\SourceLeaf{675}
\ArithmeticHead{40}{L}{H. GEPPERT, ÜBER GAUSS’ ARBEITEN ZUR MECHANIK UND POTENTIALTHEORIE.}
\makeatletter\tagsleft@true\let\veqno\@@leqno\makeatother
\renewcommand{\thefootnote}{\arabic{footnote})}\setcounter{footnote}{71}
\begin{center}\textbf{IV. Potentialtheorie.}\par\medskip
\textls[100]{1. Die »Allgemeinen Lehrsätze«.}\end{center}
\noindent\hspace*{\parindent}Die in der \textit{Theoria attractionis} entwickelten Methoden hat \textsc{Gauss} in der
Folgezeit hauptsächlich für die Potentialtheorie fruchtbar gemacht, die er im
Anschluss an seine funktionentheoretischen und erdmagnetischen Arbeiten bearbeitete.
Seine hierauf bezüglichen Untersuchungen hat er in der grossen Abhandlung
\textit{Allgemeine Lehrsätze in Beziehung auf die im verkehrten Verhältnisse des Quadrats
der Entfernung wirkenden Anziehungs- und Abstossungs-Kräfte} niedergelegt,
die in den »Resultaten aus den Beobachtungen des magnetischen Vereins« im
Jahre 1839 erschien\footnote{Werke V, S. 195—242. \textsc{Ostwalds} Klassiker, Band 2, 1889.}. Wir wollen erst den Inhalt dieser Abhandlung, soweit
er vom mechanischen und analytischen Gesichtspunkt aus interessiert,
besprechen und dann auf die Entstehung der \textsc{Gauss}schen Gedankengänge eingehen.

\textsc{Lagrange} hat, wie wir schon oben erwähnten, als erster die Zweckmässigkeit
der Einführung des Potentials erkannt, und \textsc{Laplace} verwandte es bei
der Behandlung des Ellipsoidproblems in seiner \textit{Mécanique céleste.} Aber erst
\textsc{Gauss} hat das Potential zum eigentlichen Gegenstand einer weittragenden
Theorie erhoben. Auch der Name »Potential«, der sich an die scholastische
Unterscheidung von »potentia« und »actus« anlehnt\footnote{Vergl. \textsc{Schaefer}, \textit{Über Gauss’ physikalische Arbeiten}, Werke XI, 2, S. 95.}, rührt von \textsc{Gauss} her.
Zwar geht zeitlich der \textsc{Gauss}schen Abhandlung die von \textsc{Green}\footnote{\textsc{George Green}, \textit{Essay on the application of mathematical analysis on the theories of electricity
and magnetism.} Nottingham 1828; abgedruckt in \textsc{Crelles} Journal f. d. reine u. angew. Mathematik 39,
44, 47, 1850—54. \textsc{Ostwalds} Klassiker, Band 61, 1895.} voran, aber
dessen Arbeiten wurden erst viel später bekannt, und die eigentliche Verbreitung
des Potentialbegriffs ist nur der \textsc{Gauss}schen Abhandlung zu danken.
Die völlige Unabhängigkeit von \textsc{Gauss} und \textsc{Green} ist sicher\footnote{\textsc{Schaefer}, a.\,a.\,O. S. 95—103.}.

Die \textsc{Gauss}sche Abhandlung besteht aus zwei Teilen; im ersten wird zum
ersten Male eine exakte Theorie des Körper- und Flächenpotentials, seiner charakteristischen
Eigenschaften und Wertverteilung gegeben; der zweite Teil
\clearpage\SourceLeaf{676}\ArithmeticHead{41}{R}{IV. POTENTIALTHEORIE.}
\HistoricalFoot{X 2 Abh. 7.}{6}
\noindent enthält den Existenz- und Eindeutigkeitsbeweis des ersten Randwertproblems
der Potentialtheorie.

Im ersten Abschnitt behandelt \textsc{Gauss} das Punkt-, Körper- und Flächenpotential.
Der Werteverlauf des Potentials wird anschaulich dargestellt durch
die Niveauflächen \(V=\mathrm{const.}\) und deren orthogonale Trajektorien, die Kraftlinien.
Ist \(p(x,y,z)\) die resultierende Kraft und \(C\) eine beliebige Kurve der
Bogenlänge \(s\), so ist
\[
p\cos(p,s)=\frac{\partial V}{\partial s},\tag{3}
\]
\noindent und daher das Linienintegral \(\int\limits_C p\cos(p,s)\,ds\) vom Wege unabhängig. Ist nur
eine Masse im Punkte \(O\) vorhanden, und \(r\) der Abstand des variablen Kurvenpunktes
von \(O\), so nimmt der letzte Satz — auch beim logarithmischen Potential — die Form an
\[
\oint\frac{\cos(r,s)}{r^2}\,ds=0,\tag{4}
\]
\noindent und erscheint damit als unmittelbare Übertragung des Satzes von der scheinbaren
Grösse aus der \textit{Theoria attractionis.} Diese von der Elektrizitätslehre
nahegelegten Begriffsbildungen kommen der Sache nach auch bei \textsc{Faraday} vor
(1831), sind aber, wie wir sehen werden, bei \textsc{Gauss} viel früher, etwa auf 1814,
zu datieren. Ausserhalb der anziehenden Massen existieren \(V\) und seine sämtlichen
Ableitungen; schon \textsc{Laplace} hatte bemerkt, dass die zweiten Ableitungen
ausserhalb der Punktmassen die Gleichung
\[
\Delta V=0\tag{5}
\]
\noindent erfüllen. 1782 war er bei der Behandlung des Ellipsoidproblems auf sie gestossen
und hatte sie zuerst in Polarkoordinaten, dann 1789 in rechtwinkligen
Koordinaten angegeben. \textsc{Gauss} kannte die Potentialgleichung unzweifelhaft aus
diesen Arbeiten, sie schlug ihm eine wesentliche Brücke zu seinen frühzeitigen
Studien über Funktionentheorie. In seiner Abhandlung wendet sich nun \textsc{Gauss}
der Übertragung dieser Potentialgleichung auf die Körperpotentiale zu, die
schon vor ihm \textsc{Poisson} untersucht hatte. Ist \(B\) der von der geschlossenen
Oberfläche \(O\) begrenzte anziehende Körper, dessen Massendichte \(f(\xi,\eta,\zeta)\) sei,
so handelt es sich um die Untersuchung der Grösse
\[
V=\iiint\limits_B\frac{f(\xi,\eta,\zeta)}{r}\,d\xi\,d\eta\,d\zeta.\tag{6}
\]
\clearpage\SourceLeaf{677}\ArithmeticHead{42}{L}{H. GEPPERT, ÜBER GAUSS’ ARBEITEN ZUR MECHANIK UND POTENTIALTHEORIE.}
\noindent Die Benutzung von Polarkoordinaten ergibt dann, dass \(V\) und seine ersten Ableitungen
sich endlich und stetig verhalten, auch wenn der Aufpunkt \(x,y,z\)
aus dem Äusseren ins Körperinnere übergeht. Dies gilt dagegen nicht mehr
für die höheren Derivierten.

Um diese zu berechnen und die \textsc{Laplace}sche Gleichung auf das Innere
von \(B\) zu übertragen, benutzt \textsc{Gauss} eine ihm eigentümliche Methode, die
wiederum auf der Anwendung von Integralsätzen beruht. Ist \(r\) der Abstand
des Aufpunktes \(x,y,z\) vom variablen Punkte \(\xi,\eta,\zeta\), so gibt die direkte Bildung
der zweiten Ableitung aus (6)
\begingroup\small
\[
\frac{\partial^2V}{\partial x^2}
=\iiint\limits_B\frac{\partial f}{\partial\xi}\,\frac{\partial r}{\partial\xi}\,\frac1{r^2}\,d\xi\,d\eta\,d\zeta
-\iint\limits_O f(\xi,\eta,\zeta)\frac{\partial r}{\partial\xi}\,\frac1r\cos(n,\xi)\,dO,\tag{7}
\]
\endgroup
\noindent somit
\[
\Delta V=\iiint\limits_B\frac{\partial f}{\partial r}\,\frac1{r^2}\,d\xi\,d\eta\,d\zeta
-\iint\limits_O f(\xi,\eta,\zeta)\frac1{r^2}\cos(r,n)\,dO,\tag{8}
\]
\noindent und diese Beziehung gilt immer, ausgenommen wenn der Aufpunkt in \(O\) fällt,
in welchem Falle die zweiten Derivierten von \(V\) dort unstetig sind, also \(\Delta V\)
eines bestimmten Sinnes entbehrt. Nun folgt bei \textsc{Gauss} der wichtige Satz,
dass
\begingroup\small
\[
\iiint\limits_B\frac{\partial f}{\partial r}\,\frac1{r^2}\,d\xi\,d\eta\,d\zeta
=\iint\limits_O\frac{f(\xi,\eta,\zeta)\cos(r,n)}{r^2}\,dO
+\left\{\begin{array}{l}0\\-4\pi f(x,y,z)\end{array}\right.\tag{9}
\]
\endgroup
\noindent ist, jenachdem der Aufpunkt ausserhalb oder innerhalb \(B\) fällt, woraus sofort
die \textsc{Laplace-Poisson}sche Gleichung
\[
\Delta V=\left\{\begin{array}{l}0\\-4\pi f(x,y,z)\end{array}\right.\tag{10}
\]
\noindent entspringt. Der Beweis von (9) wird durch Zerlegung in Elementarkegel mit
der Spitze im Aufpunkt geführt, und überhaupt tritt dieser Satz den drei
letzten Theoremen der \textit{Theoria attractionis} an die Seite, aus denen er sich
sofort ableiten lässt.

Durch diesen Beweis füllte \textsc{Gauss} eine empfindliche Lücke aus, die
\textsc{Poisson}, von dem die Gleichung (10) herrührt, gelassen hatte. \textsc{Poisson}\footnote{\textsc{Poisson}, \textit{Remarques sur une équation qui se présente dans la théorie des attractions des sphéroïdes.}
Nouv. Bull. de la soc. philomatique 3, 1813, S. 388—392.} hatte
\clearpage\SourceLeaf{678}\ArithmeticHead{43}{R}{IV. POTENTIALTHEORIE.}
\HistoricalFoot{}{6\textsuperscript{*}}
\noindent 1813 bei Gelegenheit des Sphäroidproblems zuerst die allgemeine Gleichung
(10) entdeckt, und für sie in der Folgezeit drei Beweise entwickelt, von denen
der zweite (1823)\footnote{\textsc{Poisson}, \textit{Mémoire sur la théorie du magnétisme en mouvement.} Mém. de l’Acad. de Paris 6,
1823, S. 455—463.} dem \textsc{Gauss}schen parallel läuft und sich ebenfalls auf die
Integralsätze der \textit{Theoria attractionis} stützt; aber \textsc{Poisson}, wie seine Nachahmer,
verfallen in den Fehler, nur die rechte Seite von (8) zu untersuchen, und es
fehlt der strenge Nachweis der Identität mit \(\Delta V\), sowie die Bemerkung, dass
die letztere Grösse an einer Unstetigkeitsstelle von \(f\) keinen eindeutigen Sinn
besitzt; es hat daher auch keinen Sinn, wenn \textsc{Poisson} im Falle des in \(O\) liegenden
Aufpunktes die rechte Seite von (10) durch \(-2\pi f(x,y,z)\) ersetzt.

Die heute klassischen Beweise der Gleichung (10) von \textsc{Dirichlet}\footnote{\textsc{P. G. Lejeune-Dirichlet}, \textit{Vorlesungen über die im umgekehrten Verhältnis des Quadrats der
Entfernung wirkenden Kräfte}, 2. Aufl. Leipzig 1887, S. 18—28.} und
\textsc{Riemann}\footnote{\textsc{B. Riemann}, \textit{Schwere, Elektrizität und Magnetismus.} Herausg. von \textsc{Hattendorff}, Hannover
1876, S. 44 ff.} stützen sich ebenfalls auf die \textsc{Gauss}schen Sätze von 1813, und
wenn \textsc{Gauss} im Art. 1 der \textit{Theoria attractionis} auf weitere Anwendungen seiner
Integralsätze hinweist, so dürfte er in erster Linie diese potentialtheoretischen
Folgerungen im Auge gehabt haben.

Nun wendet sich \textsc{Gauss} der Analyse der charakteristischen Eigenschaften
des Flächenpotentials zu. Dieser Begriff war durch die Arbeiten von \textsc{Coulomb}
und \textsc{Poisson} zur Elektrizitätslehre begründet worden, erscheint aber erst
bei \textsc{Gauss} als Träger einer selbständigen Theorie. \textsc{Poisson} fasst das Flächenpotential
immer als ein von sehr dünnen Massenschichten erzeugtes Körperpotential
auf, während \textsc{Gauss} ideale, nur flächenhaft ausgebreitete Massen
einführt. \textsc{Poisson}\footnote{\textsc{Poisson}, \textit{Mémoire sur la distribution de l’électricité à la surface des corps conducteurs.} Mém.
de l’Institut, 12, Paris 1811.} hatte 1811 auf synthetischem Wege gezeigt, dass beim
Überschreiten der anziehenden Fläche ein solches Potential der Beziehung
genügt
\[
\frac{\partial V}{\partial n_+}-\frac{\partial V}{\partial n_-}=-4\pi f,\tag{11}
\]
\noindent worin \(n\) die innere Normalenrichtung, \(f\) die Flächendichte bezeichnet; \textsc{Cauchy}\footnote{\textsc{A. L. Cauchy}, \textit{De la différence entre les attractions.\, .\, .} Bull. de la soc. philomatique 1815,
S. 53—56.}
\clearpage\SourceLeaf{679}\ArithmeticHead{44}{L}{H. GEPPERT, ÜBER GAUSS’ ARBEITEN ZUR MECHANIK UND POTENTIALTHEORIE.}
\noindent lieferte 1815 den ersten analytischen Beweis dieser Formel, jedoch bleibt
\textsc{Gauss} und \textsc{Green} das Verdienst einer ersten strengen und systematischen
Untersuchung der Flächenpotentiale und ihrer Derivierten.

Das Flächenpotential \(V\) ist überall endlich und stetig, seine Derivierten
dagegen nur bis auf die Punkte der Fläche selbst. Ist \(A\) ein Punkt der
letzteren, \(t\) eine beliebige hindurchgehende Richtung und \(p_t^A\) die in \(A\) in dieser
Richtung wirkende Kraftkomponente, so beweist \textsc{Gauss} die folgende Formel:
\[
\left.\frac{\partial V}{\partial t}\right|_A=p_t^A\mp2\pi f^A\cos(n,t),\tag{12}
\]
\noindent jenachdem \(dt\) positiv oder negativ ist. Aus ihr folgt speziell die \textsc{Poisson}sche
Formel (11), (12) geht aber insofern weiter, als sie erstens die Annäherung an
\(A\) in beliebiger Richtung zulässt, und zweitens nicht nur die Werte des Potentials
auf den beiden Seiten der Fläche mit einander, sondern diese wiederum
mit der Kraft in der Fläche selbst verbindet. Es findet also beim Überschreiten
der Fläche, wie \textsc{Gauss} sich ausdrückt, ein doppelter Sprung statt.
An den auf exakten Integralsätzen beruhenden \textsc{Gauss}schen Beweis von (12)
schliessen sich auch die späteren Beweise von \textsc{Dirichlet}, \textsc{Riemann} u.\,a. an.

Mit dem Art. 19 beginnt die \textsc{Gauss} eigentümliche Theorie, deren Grundlage
die Reziprozität zwischen Attraktionszentrum und Aufpunkt bildet, d.\,h.
die Bemerkung, dass, wenn in \(P_i\) die Massen \(M_i\), in \(p_\kappa\) die Massen \(m_\kappa\) sitzen
und die hervorgerufenen Potentiale \(V\) bezw. \(v\) sind, dann die Gleichung gilt
\[
\sum_i M_i v(P_i)=\sum_\kappa m_\kappa V(p_\kappa),\tag{13}
\]
\noindent die auch auf Körper- und Flächenpotentiale sinngemäss übertragen werden
kann. Nimmt man etwa für \(v\) das Flächenpotential einer Kugelfläche \(O\) vom
Radius \(R\), so folgt aus (13) die \textsc{Gauss}sche Fassung des Satzes vom arithmetischen
Mittel:
\begingroup\small
\[
\text{Mittelwert von }V\text{ über die Kugel}=\frac1{4\pi R^2}\iint V\,dO=V^0+\frac{M^0}{R};\tag{14}
\]
\endgroup
\noindent worin \(V^0\) das Potential der ausserhalb \(O\) befindlichen Massen im Kugelmittelpunkt,
\(M^0\) die im Kugelinnern befindliche Gesamtmasse bezeichnen. Dieser
Satz, der später von \textsc{C. Neumann} für eine beliebige geschlossene Fläche formuliert
wurde, gestattet eine grosse Anzahl von Folgerungen, z.\,B. folgt aus
ihm der Fundamentalsatz der analytischen Fortsetzung: In einem massenfreien,
\clearpage\SourceLeaf{680}\ArithmeticHead{45}{R}{IV. POTENTIALTHEORIE.}
\noindent zusammenhängenden Raume ist das Potential konstant, sobald es in einem
Teile dieses Raumes konstant ist.

Der Satz von der scheinbaren Grösse, d.\,h. der Satz 4 der \textit{Theoria attractionis},
ergibt als direkte potentialtheoretische Folgerung die Formel
\[
\text{Kraftfluss durch }O=\iint\limits_O p_n\,dO=4\pi M+2\pi M';\tag{15}
\]
\noindent darin bedeutet \(O\) eine geschlossene Fläche, \(M\) die innerhalb, \(M'\) die auf \(O\)
liegende Gesamtmasse und \(p_n\) die erzeugte Kraftkomponente in Richtung der
inneren Normalen. Die Verbindung mit dem Diskontinuitätssatz (12) ergibt
den sog. \textsc{Gauss}schen Satz oder zweiten Mittelwertsatz
\[
\iint\limits_O\frac{\partial V}{\partial n}\,dO=4\pi M,\tag{16}
\]
\noindent worin \(M\) die von \(O\) eingeschlossene Gesamtmasse ist, zu der \(M'\) hinzuzunehmen
ist oder nicht, jenachdem man \(\frac{\partial V}{\partial n_-}\) oder \(\frac{\partial V}{\partial n_+}\) betrachtet. Die Auffindung der
drei letztgenannten Sätze, die heutigentages aus der allgemeinen \textsc{Green}schen
Formel abgeleitet zu werden pflegen, reicht bei \textsc{Gauss} sehr weit, jedenfalls
in das erste Jahrzehnt des 19. Jahrhunderts zurück.

Nun wendet sich \textsc{Gauss} den Wertverteilungssätzen für \(V\) zu. Den Ausgangspunkt
bietet ihm die durch Prismenzerlegung des Raumes gewonnene
Formel
\begingroup\small
\[
\iiint\limits_B\left\{\left(\frac{\partial V}{\partial x}\right)^2+\left(\frac{\partial V}{\partial y}\right)^2+\left(\frac{\partial V}{\partial z}\right)^2\right\}\,dx\,dy\,dz
+\iint\limits_O V\frac{\partial V}{\partial n}\,dO=0,\tag{17}
\]
\endgroup
\noindent die ganz in der \textsc{Green}schen Manier abgeleitet wird und gleichfalls ein Spezialfall
des allgemeinen \textsc{Green}schen Satzes ist. Vor allem ergibt sich aus ihr der
Eindeutigkeitssatz der Potential- und Funktionentheorie: Liegen innerhalb und
auf einer geschlossenen Niveaufläche von \(V\) keine Massen, so ist \(V\) im Innern
derselben konstant. \textsc{Gauss} hebt diesen Satz in seiner Selbstanzeige besonders
hervor und präzisiert damit eine von \textsc{Poisson} gemachte Bemerkung\footnote{\textsc{Schaefer}, a.\,a.\,O. S. 99.}. Als weiterer
Satz dieser Gruppe ist der folgende von Wichtigkeit: Ist \(O\) eine geschlossene
Niveaufläche \(V=A\), und sind die Massen nur innerhalb und auf \(O\) gelegen, so
ist ausserhalb \(O\): \(|V|<|A|\) und \(\operatorname{sign}V=\operatorname{sign}A\), bezw. \(V=0\), wenn \(A=0\)
\clearpage\SourceLeaf{681}\ArithmeticHead{46}{L}{H. GEPPERT, ÜBER GAUSS’ ARBEITEN ZUR MECHANIK UND POTENTIALTHEORIE.}
\noindent ist, und diese Fälle treten ein, jenachdem \(M\neq0\) oder \(M=0\) ist. \textsc{Gauss} überträgt
diesen Satz auch auf nichtgeschlossene Flächen. Alle diese Wertverteilungssätze
beruhen hauptsächlich auf dem ersten Mittelwertsatz; durch
systematische Anwendung dieses Hilfsmittels hat später \textsc{C. Neumann}\footnote{\textsc{C. Neumann}, \textit{Untersuchungen über das Newtonsche und logarithmische Potential}, Leipzig 1877,
S. 27—52.} die
Untersuchung über Extremwerte des Potentials zum Abschluss gebracht und
allgemeine Sätze erhalten, in denen die \textsc{Gauss}schen Sätze als Spezialfälle enthalten
sind.

An diese Sätze schliesst nun \textsc{Gauss} die Lösung des Randwertproblems
an. Es sei \(O\) eine geschlossene Fläche; dann werden zunächst nur die zugehörigen
Flächenpotentiale \(V\) der Gesamtmasse \(M\) betrachtet, die zu einer
gleichartigen Verteilung gehören, d.\,h. einer solchen, bei der die Massendichte
\(f\) überall einerlei Zeichen besitzt. Der erste Satz, auf den sich alles Folgende
stützt, ist der folgende: Ist \(U\) eine auf \(O\) stetig vorgegebene Funktion, so
nimmt das Integral
\[
\iint\limits_O(V-2U)f\,dO=\Omega\tag{18}
\]
\noindent bei gleichartiger Verteilung \(f\geq0\) ein Minimum an, bei dem dort, wo \(f>0\)
ist, \(W=V-U\) konstant und dort, wo \(f=0\) ist, \(W\) nicht kleiner als diese
Konstante ist. Dass der Schluss auf die Existenz des betrachteten Minimums
nicht bindend ist, ist bekannt; zur Kritik vgl. den Essay von \textsc{Bolza}\footnote{\textsc{Bolza}, \textit{Gauss und die Variationsrechnung}, Werke X, 2, S. 77—84.}.

Aus dem genannten Satze folgert nun \textsc{Gauss}, dass es eine und nur eine
gleichartige Belegung von \(O\) gibt, bei der \(\iint\limits_O Vf\,dO\) zum Minimum wird, und
zwar ist dann \(V\) auf \(O\) konstant, und in keinem Teile von \(O\) ist \(f=0\).
Ein Kunstgriff vermittelt den Übergang zu ungleichartigen Verteilungen der
Gesamtmasse \(M\); es folgt, dass es stets eine und nur eine gleich- oder ungleichartige
Verteilung gibt, bei der auf \(O\) die Differenz \(W=V-U\) konstant
ist, und diese Konstante kann jeden Wert annehmen, indem man \(M\)
passend wählt. Dies ist der Existenz- und Eindeutigkeitssatz des ersten Randwertproblems.
Er lässt mannigfache Anwendungen zu: Nimmt man für \(U\) das
Potential eines gegebenen Massensystems, z.\,B. von Massen, die \(O\) umschliesst, so
\clearpage\SourceLeaf{682}\ArithmeticHead{47}{R}{IV. POTENTIALTHEORIE.}
\noindent kann man dessen Wirkung ausserhalb \(O\) durch die einer auf \(O\) verteilten Masse
ersetzen; ist speziell \(O\) Niveaufläche des Potentials \(U\), so lässt sich die neue Belegungsdichte
auf \(O\) vermittels (11) bestimmen zu \(f=\frac{p}{4\pi}\), wo \(p\) die von \(U\)
auf \(O\) erzeugte Kraft ist. In seiner Anzeige macht \textsc{Gauss} eine weitere Anwendung:
Eine gegebene Masse \(M\) kann auf \(O\) stets so verteilt werden, und
zwar auf eine und nur eine Art, dass im Innern von \(O\) sich die Anziehungskräfte
zerstören; dieses Theorem ist eine Verallgemeinerung des \textsc{Mac-Laurin}schen
Satzes bei Ellipsoiden.

Die genannten Sätze hat \textsc{Gauss} schon 1832 besessen, denn er wandte sie
in diesem Jahre in der \textit{Intensitas vis magneticae}\footnote{Werke V, S. 87.} und ebenso 1838 in der \textit{Allgemeinen
Theorie des Erdmagnetismus} auf die Erscheinungen des Erdmagnetismus
an. Es sei noch bemerkt, dass diese letzten \textsc{Gauss}schen Fundamentalsätze
allgemeiner sind als die entsprechenden der \textsc{Green}schen Abhandlung;
denn bei letzterem muss die Belegungsfunktion \(U\) Potentialfunktion sein,
während sie bei \textsc{Gauss} eine beliebige, stetige und endliche Funktion ist.
\clearpage\endgroup
